\begin{question}

  Evaluate the $x-p$ uncertainty product
  $\dispersion{x}\dispersion{p}$ for a one-dimensional particle
  confined between two rigid walls

  \begin{equation}
    V =
    \begin{cases} 
      0 & \text{for } 0 < x < a,\\
      \infty & \text{otherwise}
    \end{cases}
  \end{equation}  

  \begin{questionpart}
    Do this for both the ground
  \end{questionpart}
  
  \begin{questionpart}
    and excited states.
  \end{questionpart}
\end{question}

\begin{purpose}
\end{purpose}

\begin{answer}

  \begin{answerpart}{You're Grounded Mr}

    Pulling out our undergrad wavey-mechanicy quantum book, we find
    that for the infinite square potential well problem

    \begin{equation}
      \tag{griffiths:2.23}
      \inlabel{griffiths:e}
      E_n = \frac{\hbar^2k_n^2}{2m} = \frac{n^2\pi^2\hbar^2}{2ma^2}
    \end{equation}

    \begin{equation}
      \tag{griffiths:2.24}
      \inlabel{griffiths:psi}
      \psi_n(x) = \sqrt{\frac{2}{a}}\sin{\frac{n\pi}{a}x}
    \end{equation}

    So, that gives us a $\psi_n$ so let's see what happens if we just
    set up our problem.

    \begin{equation}
      \tag{1.251}
      \inlabel{exp:p}
      \begin{split}
        \mel{\beta}{p}{\alpha}
        &= \int_0^a dx' \braket{\beta}{x'}
        \left(
        -i \hbar \pdv{}{x'} \braket{x'}{a}
        \right)\\
        &= \int_0^a dx'
        \psi_{\beta}^{*}(x')
        \operatorP{x'}
        \psi_{\alpha}(x')\\
      \end{split}
    \end{equation}

    Cramming equation \inref{griffiths:e} into \inref{exp:p} we find
    
    \begin{equation}
      \inlabel{expp}
      \begin{split}
        \mel{\psi_n}{\hat{p}}{\psi_n}
        &=
        \int_0^a dx'
        \psi_{n}^{*}(x')
        \operatorP{x'}
        \psi_{n}(x')\\
        &=
        \int_0^a dx'
        \sqrt{\frac{2}{a}}\sin{\frac{n\pi}{a}x'}
        \operatorP{x'}
        \sqrt{\frac{2}{a}}\sin{\frac{n\pi}{a}x'}
        \\
        &=
        -i \hbar \frac{2}{a}
        \int_0^a dx'
        \sin{\frac{n\pi}{a}x'}
        \frac{n\pi}{a}
        \cos{\frac{n\pi}{a}x'}
        \\
        &=
        -i \hbar \frac{2}{a}
        \frac{n\pi}{a}
        \int_0^a dx'
        \sin{\frac{n\pi}{a}x'}
        \cos{\frac{n\pi}{a}x'}
        \\
        &=
        -i \hbar
        \frac{n\pi}{a^2}
        \int_0^a dx'
        \sin{\frac{2n\pi}{a}x'}
        \\
        &=
        i \hbar
        \frac{n\pi}{a^2}
        \frac{a}{2n\pi}
        \eval{
          \cos{\frac{2n\pi}{a}x'}
          }_0^a
        \\
        &=
        i \hbar
        \frac{1}{2a}
        \left(
        \cos{\frac{2n\pi}{a}a}
        -
        \cos{\frac{2n\pi}{a}0}
        \right)
        \\
        &=
        i \hbar
        \frac{1}{2a}
        \left(
        \cos{2n\pi}
        -
        \cos{0}
        \right)
        \\
        &= 0
        \\
      \end{split}
    \end{equation}

    Before we go much further, let's apply a little bit of physical
    reasoning about this one.  As usual, we're going to consider the
    physics of the problem \emph{after} we go through the pain.  When
    we look at the wave equation, we see that it is symmetrical about
    $a/2$.  So we're just as likely to find the particle on the left
    as we are on the right.  So, in theory, the expectation-value of
    momentum for any symmetrical wave function should be 0 just like
    the expectation-value for position should be the point of symmetry
    ($a/2$ in this case), though this is not something we have proven.

    \begin{equation}
      \begin{split}
        \mel{\psi_n}{\hat{p}^2}{\psi_n}
        &=
        \int_0^a dx'
        \psi_{n}^{*}(x')
        \left(
        (-1)^2i^2\hbar^2\frac{\partial^2}{\partial x^{\prime 2}}
        \right)
        \psi_{n}(x')\\
        &=
        \int_0^a dx'
        \sqrt{\frac{2}{a}}\sin{\frac{n\pi}{a}x'}
        \left(
        (-1)^2i^2\hbar^2\frac{\partial^2}{\partial x^{\prime 2}}
        \right)
        \sqrt{\frac{2}{a}}\sin{\frac{n\pi}{a}x'}
        \\
        &=
        -\hbar^2
        \frac{2}{a}
        \int_0^a dx'
        \sin{\frac{n\pi}{a}x'}
        \left(
        \frac{\partial^2}{\partial x^{\prime 2}}
        \right)
        \sin{\frac{n\pi}{a}x'}
        \\
        &=
        -\hbar^2
        \frac{2}{a}
        \frac{n\pi}{a}
        \int_0^a dx'
        \sin{\frac{n\pi}{a}x'}
        \left(
        \pdv{}{x'}
        \right)
        \cos{\frac{n\pi}{a}x'}
        \\
        &=
        \hbar^2
        \frac{2}{a}
        \frac{n^2\pi^2}{a^2}
        \int_0^a dx'
        \sin{\frac{n\pi}{a}x'}
        \sin{\frac{n\pi}{a}x'}
        \\
        &=
        \hbar^2
        \frac{2}{a}
        \frac{n^2\pi^2}{a^2}
        \int_0^a dx'
        \sin^2{\frac{n\pi}{a}x'}
        \\
      \end{split}
    \end{equation}

    Recalling our trig identity from equation \qref{trigCos2A}

    \begin{equation}
      \begin{alignedat}{2}
        & & 
        \cos{2\alpha}
        &= 1 - 2\sin^2{\alpha}\\
        \implies \quad & &
        \sin^2{\alpha} &= \frac{1 - \cos{2\alpha}}{2}\\
      \end{alignedat}
    \end{equation}

    Substituting in that trig identity, we find that

    \begin{equation}
      \inlabel{expp2}
      \begin{split}
        \mel{\psi_n}{\hat{p}^2}{\psi_n}
        &=
        \hbar^2
        \frac{2}{a}
        \frac{n^2\pi^2}{a^2}
        \int_0^a dx'
        \sin^2{\frac{n\pi}{a}x'}
        \\
        &=
        \hbar^2
        \frac{2}{a}
        \frac{n^2\pi^2}{a^2}
        \int_0^a dx'
        \frac{1 - \cos{\frac{2n\pi}{a}x'}}{2}
        \\
        &=
        \hbar^2
        \frac{2}{a}
        \frac{n^2\pi^2}{a^2}
        \left(
        \frac{a}{2}
        -
        \frac{1}{2}
        \int_0^a dx'
        \cos{\frac{2n\pi}{a}x'}
        \right)
        \\
        &=
        \hbar^2
        \frac{2}{a}
        \frac{n^2\pi^2}{a^2}
        \left(
        \frac{a}{2}
        -
        \frac{1}{2}
        \eval{
          \frac{a}{2n\pi}\sin{\frac{2n\pi}{a}x'}
        }_0^a
        \right)
        \\
        &=
        \hbar^2
        \frac{2}{a}
        \frac{n^2\pi^2}{a^2}
        \left(
        \frac{a}{2}
        \right)
        \\
        &=
        \hbar^2
        \frac{n^2\pi^2}{a^2}
        \\
      \end{split}
    \end{equation}

    Recalling equation \inref{griffiths:e} from earlier, and comparing
    we see that:

    \begin{equation}
      \begin{split}
        E_n
        &= \frac{1}{2m}\frac{n^2\pi^2\hbar^2}{a^2}
        \\
        \mel{\psi_n}{\hat{p}^2}{\psi_n}
        &= \frac{n^2\pi^2\hbar^2}{a^2}
        \\
      \end{split}
    \end{equation}
    
    So, what we ultimately find here is that the energy of the system
    is almost pulled directly from classical mechanics with $P^2/2m$.
    This is one of those moments where you pause and say ``huh...I
    guess it actually works after all.''

    Moving on from our reveries to the expecatation of position, we
    note that as we discussed earlier, a symmetric wave function
    should have an expectation value of position equal to the point of
    symmetry.  So I'm going to put it out there and predict that the
    below equation is going to work out to $\expval{\hat{x}} = a/2$.
    Let's work it out, shall we?

    \begin{equation}
      \begin{split}
        \mel{\psi_n}{\hat{x}}{\psi_n}
        &=
        \int_0^a dx'
        \psi_{n}^{*}(x')
        \operatorP{x'}
        \psi_{n}(x')\\
        &=
        \int_0^a dx'
        \sqrt{\frac{2}{a}}\sin{\frac{n\pi}{a}x'}
        x'
        \sqrt{\frac{2}{a}}\sin{\frac{n\pi}{a}x'}
        \\
        &=
        \frac{2}{a}
        \int_0^a dx'
        x'
        \sin^2{\frac{n\pi}{a}x'}
        \\
      \end{split}
    \end{equation}

    Substituting the same trig identity from the
    $\mel{\psi_n}{\hat{p}^2}{\psi_n}$ computation in to make the
    integral easier

    \begin{equation}
      \begin{split}
        \mel{\psi_n}{\hat{x}}{\psi_n}
        &=
        \frac{2}{a}
        \int_0^a dx'
        x'
        \sin^2{\frac{n\pi}{a}x'}
        \\
        &=
        \frac{2}{a}
        \int_0^a dx'
        x'
        \left(\frac{1 - \cos{\frac{2n\pi}{a}x'}}{2}\right)
        \\
        &=
        \frac{1}{a}
        \left\{
        \int_0^a dx'
        x'
        -
        \int_0^a dx'
        x'
        \cos{\frac{2n\pi}{a}x'}
        \right\}
        \\
        &=
        \frac{1}{a}
        \left\{
        \eval{\frac{x^2}{2}}_0^a
        -
        \int_0^a dx'
        x'
        \cos{\frac{2n\pi}{a}x'}
        \right\}
        \\
        &=
        \frac{a}{2}
        -
        \frac{1}{a}
        \int_0^a dx'
        x'
        \cos{\frac{2n\pi}{a}x'}
        \\
      \end{split}
    \end{equation}

    Choosing our $u$ and $v'$ for an integration by parts so we can
    get rid of the $x'$:

    \begin{equation}
      \begin{split}
        \int u'v &= uv - \int v'u\\
        u &= \frac{a}{2n\pi}\sin{\frac{2n\pi}{a}x'}\\
        u' &= \cos{\frac{2n\pi}{a}x'}\\
        v &= x'\\
        v' &= 1 \\
      \end{split}
    \end{equation}

    \begin{equation}
      \inlabel{expx}
      \begin{split}
        \mel{\psi_n}{\hat{x}}{\psi_n}
        &=
        \frac{a}{2}
        -
        \frac{1}{a}
        \int_0^a dx'
        x'
        \cos{\frac{2n\pi}{a}x'}
        \\
        &=
        \frac{a}{2}
        -
        \frac{1}{a}
        \eval{
        \left(
        \frac{a}{2n\pi}x'\sin{\frac{2n\pi}{a}x'}
        - \frac{a}{2n\pi}\sin{\frac{2n\pi}{a}x'}
        \right)
        }_0^a
        \\
        &=
        \frac{a}{2}
        -
        \frac{1}{2n\pi}
        \eval{
        \left(
        \sin{\frac{2n\pi}{a}x'}(x' - 1)
        \right)
        }_0^a
        \\
        &=
        \frac{a}{2}
        -
        \frac{1}{2n\pi}
        \left(
        \sin{\frac{2n\pi}{\cancel{a}}\cancel{a}}(a - 1)
        - 
        \cancel{\sin{\frac{2n\pi}{a}0}(0 - 1)}
        \right)
        \\
        &=
        \frac{a}{2}
        -
        \frac{1}{2n\pi}
        \left(
        \cancel{\sin{2n\pi}(a - 1)}
        \right)
        \\
        &= \frac{a}{2}\\
      \end{split}
    \end{equation}

    So, we see that our physical reasoning actually worked out.  High
    fives all around.  Let's see what happens with $\expval{x^2}$.

    \begin{equation}
      \begin{split}
        \mel{\psi_n}{\hat{x}^2}{\psi_n}
        &=
        \int_0^a dx'
        \psi_{n}^{*}(x')
        \operatorP{x'}
        \psi_{n}(x')\\
        &=
        \int_0^a dx'
        \sqrt{\frac{2}{a}}\sin{\frac{n\pi}{a}x'}
        x^{'2}
        \sqrt{\frac{2}{a}}\sin{\frac{n\pi}{a}x'}
        \\
        &=
        \frac{2}{a}
        \int_0^a dx'
        x^{'2}
        \sin^2{\frac{n\pi}{a}x'}
        \\
      \end{split}
    \end{equation}

    Using the same trig substitution as before:

    \begin{equation}
      \begin{split}
        \mel{\psi_n}{\hat{x}^2}{\psi_n}
        &=
        \frac{2}{a}
        \int_0^a dx'
        x^{'2}
        \sin^2{\frac{n\pi}{a}x'}
        \\
        &=
        \frac{2}{a}
        \int_0^a dx'
        x^{'2}
        \left(\frac{1 - \cos{\frac{2n\pi}{a}x'}}{2}\right)
        \\
        &=
        \frac{1}{a}
        \left\{
        \int_0^a dx'
        x^{'2}
        -
        \int_0^a dx'
        x^{'2}
        \cos{\frac{2n\pi}{a}x'}
        \right\}
        \\
        &=
        \frac{1}{a}
        \left\{
        \eval{\frac{x^3}{3}}_0^a
        -
        \int_0^a dx'
        x^{'2}
        \cos{\frac{2n\pi}{a}x'}
        \right\}
        \\
        &=
        \frac{a^2}{3}
        -
        \frac{1}{a}
        \int_0^a dx'
        x^{'2}
        \cos{\frac{2n\pi}{a}x'}
        \\
      \end{split}
    \end{equation}

    We'll select the same portions for our integration by parts, but
    this time around we're going to need to integrate twice and for
    the second one we'll get to leverage our earlier result.

    \begin{equation}
      \begin{split}
        \int u'v &= uv - \int v'u\\
        u &= \frac{a}{2n\pi}\sin{\frac{2n\pi}{a}x'}\\
        u' &= \cos{\frac{2n\pi}{a}x'}\\
        v &= x^{'2}\\
        v' &= 2x' \\
      \end{split}
    \end{equation}
    
    \begin{equation}
      \begin{split}
        \mel{\psi_n}{\hat{x}^2}{\psi_n}
        &=
        \frac{a^2}{3}
        -
        \frac{1}{a}
        \int_0^a dx'
        x^{'2}
        \cos{\frac{2n\pi}{a}x'}
        \\
        &=
        \frac{a^2}{3}
        -
        \frac{1}{a}
        \left(
        uv - \int_0^a v'u
        \right)
        \\
        &=
        \frac{a^2}{3}
        -
        \frac{1}{a}
        \left(
        \eval{
          x^{'2}\frac{a}{2n\pi}\sin{\frac{2n\pi}{a}x'}
        }_0^a
        - \int_0^a
        2x'\frac{a}{2n\pi}\sin{\frac{2n\pi}{a}x'}
        \right)
        \\
        &=
        \frac{a^2}{3}
        -
        \frac{a^2}{2n\pi}\sin{\frac{2n\pi}{\cancel{a}}\cancel{a}}
        +
        \frac{1}{\cancel{a}}
        \int_0^a
        \cancel{2}x'\frac{\cancel{a}}{\cancel{2}n\pi}\sin{\frac{2n\pi}{a}x'}
        \\
        &=
        \frac{a^2}{3}
        -
        \frac{a^2}{2n\pi}\cancel{\sin{2n\pi}}
        +
        \frac{1}{n\pi}
        \int_0^a
        x'\sin{\frac{2n\pi}{a}x'}
        \\
        &=
        \frac{a^2}{3}
        +
        \frac{1}{n\pi}
        \int_0^a
        x'\sin{\frac{2n\pi}{a}x'}
        \\
      \end{split}
    \end{equation}

    Repeating our integration by parts with new values

    \begin{equation}
      \begin{split}
        \int u'v &= uv - \int v'u\\
        u &= -\frac{a}{2n\pi}\cos{\frac{2n\pi}{a}x'}\\
        u' &= \sin{\frac{2n\pi}{a}x'}\\
        v &= x'\\
        v' &= 1 \\
      \end{split}
    \end{equation}
    
    
    We'll repeat our integration by parts, but noting that it's now
    $x'$ instead of $x^{'2}$

    \begin{equation}
      \inlabel{expx2}
      \begin{split}
        \mel{\psi_n}{\hat{x}^2}{\psi_n}
        &=
        \frac{a^2}{3}
        +
        \frac{1}{n\pi}
        \int_0^a
        x'\sin{\frac{2n\pi}{a}x'}
        \\
        &=
        \frac{a^2}{3}
        +
        \frac{1}{n\pi}
        \left(
        uv - \int_0^a v'u
        \right)
        \\
        &=
        \frac{a^2}{3}
        +
        \frac{1}{n\pi}
        \left(
        -\eval{
          x'\frac{a}{2n\pi}\cos{\frac{2n\pi}{a}x'}
        }_0^a
        -
        \int_0^a
        -\frac{a}{2n\pi}\cos{\frac{2n\pi}{a}x'}
        \right)
        \\
        &=
        \frac{a^2}{3}
        +
        \frac{1}{n\pi}
        \left(
        -\frac{a^2}{2n\pi}\cos{\frac{2n\pi}{\cancel{a}}\cancel{a}}
        +
        \frac{a}{2n\pi}
        \eval{
          \left[
          \frac{a}{2n\pi}\sin{\frac{2n\pi}{a}x'}
          \right]
        }_0^a
        \right)
        \\
        &=
        \frac{a^2}{3}
        +
        \frac{1}{n\pi}
        \left(
        -\frac{a^2}{2n\pi}
        +
        \frac{a^2}{4n^2\pi^2}
        \sin{\frac{2n\pi}{\cancel{a}}\cancel{a}}
        \right)
        \\
        &=
        \frac{a^2}{3}
        -
        \frac{1}{n\pi}
        \frac{a^2}{2n\pi}
        \\
        &=
        \frac{a^2}{3} - \frac{a^2}{2n^2\pi^2}
        \\
      \end{split}
    \end{equation}

    
    Putting it all together, we find that:

    \begin{equation}
      \inlabel{handy}
      \begin{alignedat}{2}
        \expval{x}^2 &= \frac{a^2}{4} & \quad{} &\text{(eq \inref{expx})}\\
        \expval{x^2} &= \frac{a^2}{3} - \frac{a^2}{2n^2\pi^2} & &\text{(eq \inref{expx2})}\\
        \expval{p}^2 &= 0 & &\text{(eq \inref{expp})}\\
        \expval{p^2} &= \frac{n^2\pi^2\hbar^2}{a^2} & &\text{(eq \inref{expp2})}\\
      \end{alignedat}
    \end{equation}

    recalling our definition of the variance

    \begin{equation*}
      \tag{1.144}
      \dispersion{A} = \expval{A^2} - \expval{A}^2
    \end{equation*}

    \begin{equation}
      \begin{split}
        \dispersion{x}\dispersion{p}
        &=
        \left(\expval{\hat{x}^2} - \expval{\hat{x}}^2\right)
        \left(\expval{\hat{p}^2} - \expval{\hat{p}}^2\right)
        \\
        &=
        \left(\frac{a^2}{3} - \frac{a^2}{2n^2\pi^2} - \frac{a^2}{4}\right)
        \left(\frac{n^2\pi^2\hbar^2}{a^2} - 0\right)
        \\
        &=
        \cancel{a^2}
        \left(\frac{1}{3} - \frac{1}{2n^2\pi^2} - \frac{1}{4}\right)
        \frac{1}{\cancel{a^2}}
        \left(n^2\pi^2\hbar^2\right)
        \\
        &=
        n^2\pi^2\hbar^2
        \left(\frac{1}{3} - \frac{1}{2n^2\pi^2} - \frac{1}{4}\right)
        \\
      \end{split}
    \end{equation}

    Let's actually pause here to note that the uncertainty is
    \emph{not} proportional in any way to $a$.  Iiiiinteresting.

    \begin{equation}
      \begin{split}
        \dispersion{x}\dispersion{p}
        &=
        n^2\pi^2\hbar^2
        \left(\frac{1}{3} - \frac{1}{2n^2\pi^2} - \frac{1}{4}\right)
        \\
        &=
        \frac{n^2\pi^2\hbar^2}{3}
        - \frac{\cancel{n^2\pi^2}\hbar^2}{2\cancel{n^2\pi^2}}
        - \frac{n^2\pi^2\hbar^2}{4}
        \\
        &=
        \frac{n^2\pi^2\hbar^2}{3}
        - \frac{n^2\pi^2\hbar^2}{4}
        - \frac{\hbar^2}{2}
        \\
        &=
        \hbar^2
        \left(
        \frac{n^2\pi^2}{12}
        - \frac{1}{2}
        \right)
        \\
      \end{split}
    \end{equation}

    So, it makes some amount of intuitive sense that the uncertainty
    isn't going to be a function of the size of the box.  Otherwise,
    consider what happens when you change the size of the box by
    changing the units.  Instead of $a$mm make it $a$cm.  You've
    changed the size of the box by a factor of 10, but otherwise
    everything remains the same.  To keep the factor of $a$ in there
    would also invite the final result to have an extra unit of
    distance in there which would preclude us from having $\hbar$
    (because then the units would be all wrong).  There's probably a
    proof hanging around about this observation, but it sure ain't in
    chapter 1 of Sakurai.

    So, we were supposed to validate our relation, so let's try it out
    with $n=0$

    \begin{equation}
      \begin{split}
        \eval{\dispersion{x}\dispersion{p}}_{n=1}
        &=
        \eval{
          \hbar^2
          \left(
          \frac{n^2\pi^2}{12}
          - \frac{1}{2}
          \right)
        }_{n=1}
        \\
        &=
        \hbar^2
        \left(
        \frac{\pi^2}{12}
        - \frac{1}{2}
        \right)
        \\
        &\approx
        \hbar^2
        \frac{1}{3}
        \\
        &>
        \hbar^2
        \frac{1}{4}
        \\
      \end{split}
    \end{equation}

    As we can see, it passes...by the skin of its teeth...but it
    passes.

  \end{answerpart}

  \begin{answerpart}{I'm So Ex-cited!}
    When we evaluate at higher energy levels, that $n^2$ term starts
    to quickly dominate as shown in figure \inref{figure:splosions}.

    \begin{figure}
      \begin{center}
        \caption{As you can see, the ground state just barely clears
          the Heisenberg limit, but it does clear it.  Pretty quickly
          thereafter the energy level's $n^2$ term dominates.}
        \includegraphics[width=0.8\linewidth]{energy_levels.pdf}
        \inlabel{figure:splosions}
      \end{center}
    \end{figure}

    This has a feeling of getting into information theory and entropy.
    The higher energy levels carry with them higher amounts of
    uncertainty...at least in this example.  Let's see what happens as
    we continue our journey.
  \end{answerpart}

\end{answer}
